\subsection{Two components of the four-node locus for \texorpdfstring{$X_9$}{X9}}\label{ssec:x9components}

The reached-component restriction in Proposition~\ref{prop:componentwise} is necessary.

Use the smooth projective toric fourfold $P=X_{\mathcal T}$ specified
by the $X_9$ triangulation of Section~\ref{ssec:x9} and Appendix~\ref{app:data-x9}. Its fan has $161$ rays
and $813$ maximal cones; this is the smooth triangulation used for
the example in the current manuscript. Put
\[
 V=H^0(P,-K_P)\setminus\{0\}.
\]
Let $S_4\subset V$ be the locus where $Y_f$ has exactly four nodes
and no other singularities.
The hypersurfaces in this locus are required to be reduced.
By Lemmas~\ref{lem:component-rank} and~\ref{lem:component-smoothness} and the local-system argument of Proposition~\ref{prop:componentwise}, $S_4$
is smooth, its local codimension is the rank of the vanishing-cycle
map, and that rank is locally constant along $S_4$.
The near-limit construction gives a member for which the rank is three
and the marked integral relation lattice is $\Z(1,-1,-1,1)$.

\begin{proposition}\label{prop:x9different}
The locus $S_4$ has at least two connected components. One has
codimension three and meets the near-limit construction. Another has
codimension four and does not meet that construction. For a sufficiently
small smoothing of any member of the latter component, the four
vanishing spheres are linearly independent over $\Q$. Consequently
their integral relation lattice and the preimage of ambient torsion
are both zero.
\end{proposition}

We prove the assertion by an explicit member. The Newton polytope
$\Delta_9$ has precisely the following twelve lattice points.
The first nine are its vertices. In the standard character lattice,
write $f_*=\sum_{i=0}^{11}c_i\chi^{n_i}$, with
\begin{equation}\label{eq:x9coefficients}
\begin{array}{c|r|r}
i&n_i&c_i\\\hline
0&(1,0,0,0)&1\\
1&(0,1,0,0)&1\\
2&(-1,-1,0,0)&668\\
3&(0,0,1,0)&324\\
4&(0,0,0,1)&324\\
5&(-4,-2,-1,0)&1944\\
6&(-4,-2,0,-1)&1944\\
7&(3,1,1,1)&108\\
8&(2,1,1,1)&432\\
9&(-1,0,0,0)&-21\\
10&(-2,-1,0,0)&2496\\
11&(0,0,0,0)&0
\end{array}
\end{equation}
The support assertion can be checked by testing the integer points
of $[-4,3]\times[-2,1]\times[-1,1]^2$ against the polar inequalities.
All boundary coefficients are nonzero.

For torus coordinates $t_1,t_2,t_3,t_4$, make the unimodular monomial
change
\begin{equation}\label{eq:x9coordinates}
 x=t_1,\qquad y=t_1t_2,\qquad
 z=t_1^2t_2t_3,\qquad w=t_1^2t_2t_4.
\end{equation}
Multiplying $f_*$ by the torus unit $xy=t_1^2t_2$ gives
\begin{equation}\label{eq:x9counterF}
\begin{split}
 F={}&y^2+(x^2-21)y+668x+2496\\
 &+324(z+w)+1944(z^{-1}+w^{-1})+108(x+4)zw.
\end{split}
\end{equation}

\subsubsection*{The four torus singularities}

The singular points of \eqref{eq:x9counterF} in the torus are the
following two pairs:
\begin{equation}\label{eq:x9nodepairs}
\begin{aligned}
 &x=5,\quad y=-2,\quad z+w=18,\quad zw=-6;\\
 &z=w=\xi,\quad \xi^2=-3,\quad
   x=-4+3\xi,\quad y=16+12\xi.
\end{aligned}
\end{equation}
Here the first pair has $z=9\pm\sqrt{87}$ and $w=9\mp\sqrt{87}$.
All four points are distinct and all coordinates are nonzero.

To prove completeness, subtract two logarithmic derivatives:
\[
 zF_z-wF_w=324(z-w)\left(1+\frac6{zw}\right).
\]
On the branch $zw=-6$, the equation reduces to
\[
 h(x,y)=y^2+(x^2-21)y+20x-96=0.
\]
The equations $h=h_x=h_y=0$ require a multiple root of the
$y$-discriminant
\[
 (x^2-21)^2-4(20x-96)=(x-5)^2(x^2+10x+33).
\]
The only multiple root is $x=5$, giving $y=-2$.
An individual logarithmic derivative then gives $z+w=18$.

On the branch $z=w$, the equations $F_z=F_y=0$ give
\[
 x=18z^{-3}-3z^{-1}-4,\qquad y=-\frac{x^2-21}{2}.
\]
After substitution, the numerators of $F_x$ and $F$, up to nonzero
rational constants, are respectively
\begin{align*}
 P_{11}={}&4z^{11}+24z^9+3z^8+4z^7-17z^6-48z^5-18z^4
            +144z^3+108z^2-216,\\
 Q_{13}={}&48z^{13}-27z^{12}+576z^{11}-18z^{10}+1712z^9+213z^8\\
          &+288z^7-576z^6-1728z^5-648z^4+3456z^3+2592z^2-3888.
\end{align*}
Their monic greatest common divisor is $z^2+3$. An exact certificate is
\[
 P_{11}=(z^2+3)P_9,\qquad Q_{13}=(z^2+3)^2Q_9,
\]
where
\begin{align*}
 P_9={}&4z^9+12z^7+3z^6-32z^5-26z^4+48z^3+60z^2-72,\\
 Q_9={}&48z^9-27z^8+288z^7+144z^6-448z^5-408z^4
          +384z^3+576z^2-432,
\end{align*}
and the nonzero resultant values are
\[
 \operatorname{Res}(P_9,Q_9)=2^{59}3^{13}13^2 29^2 31\cdot9311,
 \qquad \operatorname{Res}(P_9,z^2+3)=3^7\cdot13\cdot31.
\]
This gives exactly the second pair in \eqref{eq:x9nodepairs}.

The Hessian determinants of $F$ in the coordinates $(x,y,z,w)$ are
\begin{equation}\label{eq:x9hessians}
 \det\operatorname{Hess}F=
 \begin{cases}
  2^6 3^9\cdot87,&\text{at either point of the first pair},\\
  -2^6 3^9(21+16\xi),&\text{at the second pair, }\xi^2=-3.
 \end{cases}
\end{equation}
They are nonzero, so each singularity is an ordinary double point.

\subsubsection*{Compactification and the boundary}

It remains to exclude singular points on the toric boundary of $P$.
For a cone $\tau$ with ray generators $r_j$, let
\[
 E_0(\tau)=\{i:\langle n_i,r_j\rangle=-1\text{ for every }r_j\in\tau\}.
\]
These are the exponents surviving on the orbit of $\tau$.
If they are affinely independent, the equation and its logarithmic
tangent derivatives cannot vanish simultaneously, since all their
coefficients are nonzero. The complete cone calculation for the fixed
triangulation is
\[
\begin{array}{c|r|r|r}
\dim\tau&\text{cones}&\text{independent }E_0&\text{remaining}\\\hline
1&161&157&4\\
2&974&971&3\\
3&1626&1626&0\\
4&813&813&0
\end{array}
\]
The four relevant rays, in the accompanying data's ordering, are
\[
 r_5=(-1,2,-1,1),\quad r_6=(-1,2,1,-1),\quad
 r_8=(1,-1,-1,-1),\quad r_{11}=(0,1,-1,-1).
\]
All seven remaining cases are listed below. A normal derivative is
computed in a smooth toric chart, then restricted to the orbit; a
singleton support gives a nonzero monomial and hence a unit there.
\begin{equation}\label{eq:x9boundarytable}
\begin{array}{c|l|l}
\tau&E_0(\tau)&\text{exclusion or additional calculation}\\\hline
(5)&0,2,3,6,7&n_0\text{ outside the affine span of the others}\\
(6)&0,2,4,5,7&n_0\text{ outside the affine span of the others}\\
(8)&1,3,4,5,6,8,9,10&x=0\text{ in \eqref{eq:x9counterF}}\\
(11)&2,3,4,5,6,7,8,10&y=0\text{ in \eqref{eq:x9counterF}}\\
(5,11)&2,3,6,7&\text{normal derivative supported on }n_0\\
(6,11)&2,4,5,7&\text{normal derivative supported on }n_0\\
(8,11)&3,4,5,6,8,10&\text{normal derivative supported on }n_9
\end{array}
\end{equation}
The first two cases are excluded by taking an affine-linear combination
of the equation and tangent derivatives which retains only the
nonzero $n_0$-term. The last three normal derivatives have respective
coefficients $1,1,-21$.

For the remaining two divisor orbits, \eqref{eq:x9coordinates} identifies
the chart of $\langle r_8,r_{11}\rangle$ with
$\C_x\times\C_y\times(\C^*)^2_{z,w}$, and its anticanonical equation
is precisely \eqref{eq:x9counterF}. On either divisor orbit a singular
point would satisfy
\[
 zw=-\frac{167}{27},\qquad
 (z-w)\left(324+\frac{1944}{zw}\right)=0.
\]
The second factor in parentheses is $1620/167\ne0$; hence
$z=w$ and $z^2=-167/27$.
On $y=0$, the other normal equation and either logarithmic derivative give
\[
 x^2=21,\qquad (x+4)z=-\frac{987}{167}.
\]
The number $x$ is real and $z$ is nonzero and purely imaginary, whereas
the last right-hand side is nonzero real, a contradiction.
On $x=0$, an individual logarithmic derivative instead gives
$z=-987/668$, contradicting $z^2=-167/27$.
At $x=y=0$, the derivative $F_y=-21$ already excludes a singularity.
Thus there are no boundary singularities.

The finite checks in this section consist only of the lattice-point
list, the unimodular cone determinants, the orbit supports and their
affine ranks, and the indicated normal derivative supports. They are
reproduced from the fixed example's fan in the accompanying exact
calculation (\coderepo; \texttt{examples/verify\_x9\_components.py}). The two remaining exclusions above are over $\C$.
In particular, $Y_{f_*}$ has exactly the four nodes displayed earlier.
It is reduced: a nonreduced divisor would be singular along a
positive-dimensional subset of its support.

\subsubsection*{Smoothing rank and distinct components}

The coefficient directions indexed by $n_{10},n_2,n_3,n_4$ are the
functions $1,x,z,w$ after the trivialization used in
\eqref{eq:x9counterF}. Their evaluations at the four nodes are linearly
independent. Indeed, suppose
\[
 \alpha+\beta x+\gamma z+\delta w=0
\]
at all four points. Subtracting the values at the first pair gives
$\gamma=\delta$. The second pair then gives
$\alpha=4\beta$ and $\gamma=-3\beta/2$.
Substitution into either point of the first pair gives $-18\beta=0$.
All four coefficients therefore vanish.

\begin{proof}[Proof of Proposition~\ref{prop:x9different}]
At an ordinary double point, the differential of its local smoothing
parameter is, up to a nonzero factor, evaluation of the coefficient
variation at that point. The preceding calculation gives rank four.
Lemmas~\ref{lem:component-rank} and~\ref{lem:component-smoothness} imply
that the four vanishing classes are independent over $\Q$, and that
$S_4$ has codimension four near $f_*$. A nonzero integral combination
cannot be zero or torsion, since either would vanish after tensoring
with $\Q$.

The rank of the vanishing-cycle map is locally constant along $S_4$,
as follows from the exterior-of-the-nodes local system used in the
componentwise continuation proof. It is therefore four on the connected
component containing $f_*$. On every component meeting the near-limit
construction it is three. These components are distinct, proving all
assertions.
\end{proof}

\begin{remark}
This proves failure of coverage for the full four-node locus in this
fixed toric ambient. It does not classify all its components, nor
determine whether every component having the rank-one full-support
relation lattice meets the near-limit construction. The original
$X'_9$ remains smoothable; the counterexample shows why its mirror
relation criterion must retain the reached-component hypothesis.
\end{remark}
